\icode{O-MD-CA-L2-261113-002}
如图所示，一{\kaishu\color{purple}足够长}的水平传送带以恒定的速度\(v_0\)向右传动。将一物体轻轻放在传送带左端，传送带与零件间的动摩擦因数为 \(\mu\)。求传送带上的划痕长度。

\begin{tikzpicture}[>=stealth,font=\scriptsize]
	\begin{scope}[rotate=0]
    	\coordinate (o) at(0,0);
		\coordinate (length) at(4,0);
		\coordinate (square) at(0,.3);
		\draw[gray] (0,0) circle(.3) (length) circle(.3);
		\draw[fill=orange!60] (square) rectangle($(square)+(.3,.3)$);
		\draw ($(o)+(0,.3)$)--($(length)+(0,.3)$) arc(90:-90:.3)--($(length)+(0,-.3)$)--($(o)+(0,-.3)$) arc(270:90:.3);
		\fill (o) circle(.07) (length) circle(.07);
		\draw[->] (1.5,.15)--node[below]{\(v_0\)}(2.5,.15);
		\coordinate (o) at(.15,0.3);
		\fill (o) circle(.05);
	\end{scope}
	\begin{scope}[rotate=0,yshift=-2cm]
		\coordinate (o) at(0,0);
		\coordinate (length) at(4,0);
		\coordinate (square) at(0,.3);
		\draw[gray] (0,0) circle(.3) (length) circle(.3);
		\draw[fill=orange!20,draw=gray] (square) rectangle($(square)+(.3,.3)$);
		\draw[fill=orange!60] ($(square)+(1.7,0)$)coordinate(a) rectangle($(square)+(2,.3)$);
		\draw ($(o)+(0,.3)$)--($(length)+(0,.3)$) arc(90:-90:.3)--($(length)+(0,-.3)$)--($(o)+(0,-.3)$) arc(270:90:.3);
		\fill (o) circle(.07) (length) circle(.07);
		\coordinate (o) at(.15,0.3);
		\fill[gray!50] (o) circle(.05);
		\fill (1.8,.3) arc(180:0:.05)--(1.8,.3) (3.3,.3) arc(-180:0:.05)--(3.3,.3);
	\end{scope}
	\draw[dashed,thin,cyan] (.15,1)--(.15,-3) (1.85,-1)--(1.85,-1.7) (3.35,-1.45)--(3.35,-3);
	\draw[<->,thin,magenta] (.15,-1.2)--node[fill=white]{\(x_1\)}(1.85,-1.2);
	\draw[<->,thin,teal] (.15,-2.8)--node[fill=white]{\(x_2\)}(3.35,-2.8);
	\draw[<->] (1.85,-1.55)--node[above]{\(\Delta x\)}(3.35,-1.55);
\end{tikzpicture}

\begin{tikzpicture}[>=stealth,font=\scriptsize]
	\draw[->,blue] (0,0)node[below left]{\(O\)}--(4,0)node[below]{\(t\)};
	\draw[->,blue] (0,0)--(0,3)node[right]{\(v\)};
\end{tikzpicture}